\section{Cars Improve Faster than Transport Systems}

Replace an engine and the street stays where it was. The bridge, parking structure, school catchment, suburban subdivision, and morning schedule remain. Transport is a useful specimen of technological viscosity because its hardware is entangled with the organization of everyday life.

\subsection{A Powertrain Is One Layer}

Electric motors offer high efficiency, rapid control, and a different maintenance profile from combustion engines. Batteries move costs into materials, manufacturing, charging, thermal management, and eventual reuse or recycling. A proper comparison follows both systems through their lifecycles.

The IEA's 2026 assessment finds lifecycle emissions savings for battery electric cars under its stated comparisons, despite higher manufacturing emissions associated with batteries \cite{iea2026evlifecycle}. The size of the saving depends on vehicle, electricity, use, and lifetime. Counting the battery bill does not justify erasing the operational benefit.

The distinction that matters here is between improving a vehicle and reorganizing movement. An electric car still occupies road and parking space. A heavy battery does not invalidate electrification; a better powertrain does not settle congestion, access, or land use.

Transition costs also need a long boundary. Charging infrastructure is an added investment; fuel extraction, refining, distribution, and combustion vehicles already have investment and operating bills. Minimizing immediate change can preserve a more expensive long-run system. Engineering can compare alternative transition paths using lifecycle cost, emissions, service quality, and risk. Mechanical least action supplies no automatic verdict.

The strongest critique of automotive spectacle is therefore a precise one: the sales chart for a new powertrain cannot stand in for the performance of the transport system.

\subsection{Optimize Access, Then Movement}

People need access to work, education, care, other people, and goods. Movement supplies some of that access. Urban design, service location, remote work, and reliable local provision can reduce how much movement is necessary.

The road-flow identity is
\begin{equation}
 q=\rho v,
\end{equation}
where $q$ is vehicles per unit time, $\rho$ is vehicles per unit distance, and $v$ is space-mean speed over a consistent segment. Increasing speed need not increase throughput because safe spacing, density, and traffic dynamics change with speed.

Vehicle technology can affect these relationships: automation may change headways, vehicle sizes may change space use, and new operating patterns may change demand. The identity itself determines none of those behavioral responses.

Private cars offer flexible access, privacy, and storage. In low-density places or for complicated care and work schedules, those benefits can be decisive. The same arrangement commits households to purchasing, parking, maintaining, insuring, and often driving a rarely fully utilized asset.

Once housing, commerce, schools, and roads develop around car ownership, individual choices reinforce the structure. A household may rationally buy a car because the system has made alternatives poor. Treating the purchase as an isolated preference misses the dependency.

“Effective movement” means movement that delivers the intended access or goods. Empty repositioning, avoidable detours, searching for parking, and trips created by poor coordination consume the network without supplying the same benefit. An efficient transport policy has to see them.

\subsection{Ownership Versus Operations}

Electrification changes propulsion. Fleet operation changes how vehicles are used. These are separate axes.

A shared fleet can spread vehicle costs over more trips, centralize maintenance, and reduce some parking needs. It also incurs dispatch, cleaning, support, repositioning, and platform costs. Empty miles and demand induced by cheaper travel can increase traffic even when the fleet uses fewer vehicles than an ownership-based alternative.

Autonomous operation changes the labor bill and operating possibilities. It adds a requirement to establish reliable behavior within an operating domain, manage exceptions, maintain sensors and software, and allocate responsibility. In mixed traffic, the vehicle acts inside a social environment designed largely around human interpretation.

Waymo's October 2026 announcements illustrate the levels that need separating: expanding service in existing markets is a different activity from preparing and validating a new city \cite{waymo2026expansion,waymo2026madrid,waymo2026toronto}. Neither an impressive drive nor an announced future market alone establishes a complete service.

High utilization can reduce embodied cost per passenger trip while increasing wear, maintenance intensity, and reliance on the operator's systems. Concentration may simplify some household interfaces while creating larger dependencies on a few service providers. The three viscosity coordinates can move in different directions.

\subsection{Could the Private Car Become a Yacht?}

A provocative possibility follows. If dependable fleets provide most ordinary journeys at attractive cost, private ownership could become less necessary and more elective. Some cars would remain prized for pleasure, control, identity, collecting, or unusual journeys. In those markets, ownership would resemble a yacht more than a utility connection.

The yacht analogy describes a conditional structural shift. Its conditions are demanding: coverage, waiting time, price, accessibility, child transport, luggage, rural service, privacy, resilience, and trust. A fleet that fails them leaves substantial reasons to own a car.

This is a better proposition to test than declaring that ownership will disappear. Watch whether households with adequate service actually give up vehicles, which trips remain expensive to serve, and how the network's total vehicle-kilometers change.

Dense corridors still favor high-capacity transit. Short trips often favor walking and cycling. Freight has different loads, schedules, and terminals. Low-density trips need flexible service. No one vehicle can take over all these tasks at the same physical cost.

Resilience also needs alternatives. A software failure, severe weather event, power shortage, or operator collapse can expose the weakness of a highly concentrated service. Spare capacity and independent routes have a bill; their value appears when a preferred route fails.

The future of transport can be sharper than the future of the car. The decisive measure is reliable access per unit of whole-system cost, including space, time, energy, harm, and resilience. Some improvements arrive inside a vehicle. Others arrive when the system needs fewer vehicles to do the same useful work.
